\lesson{2}{Nov 01 2021 Mon (12:18:54)}{Public Policy}{Unit 3}

\subsubsection*{Types of Policies}

There are two major types of \textbf{Public Policy}:

\begin{itemize}
    \item Foreign
    \item Domestic
\end{itemize}

\textbf{Foreign Policy} applies to relations with other countries. \textbf{Domestic Policies} apply within the country. \textbf{Domestic Policy} and \textbf{Foreign Policy} decisions can have an affect on each other.

\textbf{Public Policies} can impact:

\begin{itemize}
    \item Local
    \item State
    \item National Operations
\end{itemize}

They depend on the legislative abilities and limitations described by the U.S. Constitution and state constitutions. Beyond the descriptions of who is affected by policies, we can also categorize them by what area of life they address.

A policy may not fit neatly into one of these categories, however. For example, a federal immigration policy may require the states to enact their own policies to enforce it, including creating economic policies to provide funding for additional law enforcement and investigators.

\paragraph{Social Welfare Policy} Social welfare policies address the well-being of citizens and include programs like education and health insurance. The goal of these policies is to \it{"Promote the General Welfare"} as described in the Preamble to the Constitution, meaning to improve citizens' quality of life. Those who benefit from social welfare programs often must meet certain age, income, or other requirements.

\paragraph{Economic Policy} Economic policies include the government's budget, which includes government taxing and spending rules. Economic policy also establishes rules for businesses. The goal of economic policy is to grow the productive capacity of the economy, while keeping unemployment and inflation low. Inflation is the general rise of prices over time. People disagree whether it is better for the economy and people when government makes rules for it.

\paragraph{Environmental Policy} Environmental policies address concerns related to the impact of human activity on Earth. The goal, of course, is to minimize the negative effects on \bf{air, water, wildlife,} and \bf{land} that result from our activities, such as clearing land for farms and driving gas-burning vehicles. Environmental policies are often controversial because they can be very expensive for American businesses.

\paragraph{Defense Policy} Defense policies include laws and programs related to maintaining the armed forces. The goal of defense policy is to protect the safety and security of the country and its citizens. Defense policy often intersects with foreign policy. Like other types of policy, it can be very controversial as people disagree on what role the United States should take in world affairs as a measure to protect its own peace and security.

\subsubsection*{Monitor Issues and Elected Officials}

Citizens in the United States have the right and responsibility to be informed about government policies that affect them. The \bf{First Amendment} protects the right of the media to report on the work of elected officials. Congress has to keep and publish a journal of its daily proceedings for public view, as do many state governments. Legislators at all levels depend on citizen input to help determine new policies or changes to public policy.

\begin{figure}[h]
    \begin{quote}
        \it{"Congress shall make no law respecting an establishment of religion, or prohibiting the free exercise thereof; or abridging the freedom of speech, or of the press; or the right of the people peaceably to assemble, and to petition the Government for a redress of grievances."}
    \end{quote}
    \caption{From the First Amendment}
\end{figure}

\subsubsection*{Policies Shaped}

Policies arise out of the lawmaking process at the:

\begin{itemize}
    \item Local
    \item State
    \item Federal
\end{itemize}

levels of government. Yet there has to be a reason for suggesting a new policy or change to an existing policy in the first place. Legislators receive \it{ideas} and \it{complaints} about all kinds of issues from:

\begin{itemize}
    \item Individuals
    \item Businesses
    \item Other Groups
\end{itemize}

They talk to other lawmakers to place the issue on the legislative agenda. Most of the states have a similar lawmaking process to the federal government, with bills starting in committees and eventually going up for vote by both halves of the legislature.

Many states grant their citizens the power of initiative or referendum, or both, which are methods through which citizens can place a policy idea on the legislative agenda. Initiative is a process that would lead to a new state policy and referendum is a process that would lead to a change or deletion of a state policy. At the local level, lawmaking varies depending on the structure of the local government according to its state constitution, yet is still dependent on people who want a solution to a particular problem.

However, legislators do not just write a bill or ordinance after talking to citizens and then have the whole lawmaking body vote on it right away. A great deal of time and analysis goes into studying an issue and shaping a policy to address it before it is ready to be enacted into law.

\subsubsection*{Costs and Benefits of Public Policy}

Every public policy has costs and benefits. The costs of a policy include:

\begin{itemize}
    \item Money
    \item Time
    \item Other Resources that go into implementing and maintaining a policy
\end{itemize}

Someone has to bear the costs, like taxpayers for example. Companies may have to hire additional workers to meet a policy requirement, which is a cost to them. The benefits are the intended effects of a policy, such as safer food or more students passing college entrance exams. Benefits could include direct payments to people, such as tuition assistance.

Public policies differ in how costs and benefits are distributed among people in a society. For example, the students in a public education program are a relatively small group of beneficiaries compared to the large group of people who help pay for the education through taxes. Some policies have distribution that is equal while others could place the cost burden on a smaller group of people. The distribution of costs and benefits can affect how citizens and elected officials feel about a policy.

Policies can also have externalities, which are unintended consequences of a decision. Externalities can be positive or negative and are called \it{"Third Party Costs"} or \it{"Third Party Benefits"}. An example of a negative externality is pollution caused by a policy encouraging new factories to operate and create jobs for a certain area.

In $1955$, after about two generations of common vehicle use, Illinois was the first state to enact a seat belt policy. Today, many states have mandatory seat belt laws.

\subsubsection*{Analyzing a Public Policy}

Changes in culture and technology happen very quickly, and it can be many years before people recognize the consequences of those changes. Public policy is the government's response to an issue or problem brought to their attention. Whether as a citizen, part of an interest group, or a legislator, people analyze public policy by learning as much as they can about the problem, considering potential solutions, and analyzing the costs and benefits of any potential solution before making a policy decision.

A problem in the United States is drivers using electronic devices, mostly cell phones, while operating a vehicle. Almost all states have banned text messaging while driving. About half have banned all cell phone use while driving. Try this activity to analyze a potential policy that would address this issue.

\newpage
